\section{Path categories, specified comparisons, and one-sign blocks}
\label{sec:paths}

All categories and indexing categories in this section are $U$-small.
The subcategory $\Sigma\subset M$ contains every object and is closed
under composition; it need not contain every isomorphism of $M$.
The nerve realization is $B$, and a weak equivalence means a member of
$W_\infty$. Mapping spaces and their equivalences are unpointed; the
empty path supplies a distinguished point when the endpoints coincide.

\begin{claim}{C.1}{Proved}{The categories and their comparison maps}
Let $\tau=\varepsilon_1\cdots\varepsilon_n$ be a finite sign word,
including the empty word. Its interval $I_\tau$ is the free category
on the oriented line with vertices $0,\ldots,n$. A type map is an
endpoint-preserving functor with nondecreasing vertex function.
Write $\mathrm{Ty}$ for their category. Put
\[
 F_\tau(a,b)=\operatorname{Fun}_\Sigma(I_\tau,M)_{a,b},
 \qquad F_\tau^{\Sigma}(a,b)\subset F_\tau(a,b).
\]
Objects send backward edges into $\Sigma$. In $F_\tau$ the vertical
components are arbitrary, with both endpoint components identities;
in $F_\tau^\Sigma$ they all belong to $\Sigma$.
The fibred Grothendieck constructions are denoted by
$\operatorname{Ch}_\Sigma(M;a,b)$ and
$\operatorname{Ch}_\Sigma^\Sigma(M;a,b)$. A morphism $(\tau,c)\to
(\sigma,d)$ is a pair $(u,\eta)$ with
$u:\tau\to\sigma$ and $\eta:c\Rightarrow d u$.
The inclusion $J$ preserves literal concatenation and empty paths.
\end{claim}
\begin{proof}
A backward edge under a type map traverses a sequence of backward edges,
possibly of length zero. Thus pullback preserves the marking condition.
Precomposition also preserves the endpoint condition and, for the
restricted version, the condition on vertical components. For composable
$(u,\eta)$ and $(v,\theta)$ the composite is
$(vu,(\theta u)\eta)$; the equations for a category follow from those
for natural transformations. Concatenation glues the final vertex of
the first interval to the initial vertex of the second, and similarly
glues functors and transformations. Its associativity and empty-word
units are strict. These formulas also prove the asserted properties of
$J$, without identifying different type maps.
\end{proof}

\begin{claim}{C.2}{Proved}{Step-by-step verification of PATHS C.P1}
For $\tau,\sigma\in\mathrm{Ty}$ the poset
$P(\tau,\sigma)$ of type maps, ordered by pointwise arrows in
$I_\sigma$, is empty or has contractible nerve. These posets form the
hom-categories of a strict monoidal $2$-category. There is a strict
contravariant action
\[
 P(\tau,\sigma)\times F_\sigma(a,b)\longrightarrow F_\tau(a,b).
\]
\end{claim}
\begin{proof}
The interval is a poset, so pointwise transformations are unique when
they exist. Induct on the number of constant-sign runs of $\sigma$.
An empty target gives a terminal mapping poset. If $\sigma$ is nonempty
and the mapping poset is nonempty, the source is nonempty. Write
$[0,a]$ and $[0,b]$ for the first source and target runs.

When their signs disagree, a map must be constant at zero on the first
source run: moving from zero through a positive number of target
edges would require the target's first sign. Deleting that source run,
while retaining its last vertex, gives an isomorphism of mapping
posets. Repeat until the signs agree; nonemptiness prevents exhaustion
of the source before a nonempty target.

When the first signs agree, put $h(j)=\max(j,b)$ and define
$(Rf)(0)=0$, $(Rf)(i)=h(f(i))$ for $i>0$. Every edge affected by $h$
lies in the first target run, where all arrows have a common direction.
At the initial source edge the new image is the arrow across that run;
on subsequent edges the image is either unchanged or lies in the same
run. Thus $Rf$ is a type map. A map cannot pass the target's first
turn within the source's first run, so that run has new values
$(0,b,\ldots,b)$. The inequalities $f\leq Rf$ in the positive case,
or $Rf\leq f$ in the negative case, are pointwise arrows and hence
natural transformations. Applying $h$ to a pointwise arrow preserves
it, so $R$ is a functor. It is idempotent and restricts to the identity
on its image. The image is exactly the full subposet with that
prescribed first run. Deleting both prescribed runs and subtracting
$b$ from the remaining target coordinates identifies it with the
mapping poset for a target having one fewer run. This proves the
contraction by induction. No naturality of the chosen contraction in
all source words is needed.

Horizontal and vertical composition preserve numerical order and
endpoint identities. Their interchange equation is the ordinary
interchange equation for natural transformations. Blockwise gluing
preserves these equations. Finally the action sends $(u,c)$ to $cu$;
on arrows $(\alpha:u\Rightarrow v,\eta:c\Rightarrow d)$ it is
\[
 cu\xrightarrow{\eta u}du\xrightarrow{d\alpha}dv.
\]
Its endpoint components are identities. Its interior components may
be arbitrary arrows of $M$. This explicitly identifies the point where
the unrestricted vertical arrows are used. For example, with
$M=[1]$, $\Sigma$ the identities and a forward interval, a nonidentity
component $0\to1$ occurs. Explicitly take $\tau=++$, $\sigma=+$,
let the two vertex maps have middle values $0$ and $1$, and pull
back the path $0\to1$. Their comparison is
$(0,0,1)\Rightarrow(0,1,1)$, with interior component $0\to1$.
Thus the corresponding action does not restrict to $F^\Sigma$.
\end{proof}

\begin{claim}{C.3}{Proved}{Descent, finality, and PATHS C.P2--C.P3}
Let $R$ be the reachability preorder of types. The diagram
$\tau\mapsto BF_\tau(a,b)$ descends coherently along
$p:\mathrm{Ty}^{\mathrm{op}}\to R^{\mathrm{op}}$.
The functor $p$ is homotopy final. Every projective sequence
$t:\mathbb N^{\mathrm{op}}\to\mathrm{Ty}$ such that each type
receives a map from some $t_n$ is homotopy final after passage to
$R$. For its actual Grothendieck construction $E_t$ and categorical
colimit $C_\infty(t)$ both maps
\[
 C_\infty(t)\longleftarrow E_t\longrightarrow
 \operatorname{Ch}_\Sigma(M;a,b)
\]
are weak equivalences. The resulting isomorphism $\beta_t$ in
$\mathrm{Hot}$ is canonical, natural under strictly commuting maps
of sequences, and transitive under the comparisons through $t*s$.
\end{claim}
\begin{proof}
Apply nerves to the action in C.2. The simplicial category with
hom-objects $NP(\tau,\sigma)$ maps to $R$, identically on objects
and by weak equivalences on every hom-object, including empty ones.
This is a Dwyer--Kan equivalence. In the projective model structures
on simplicial-set-valued enriched diagrams, all these hom-objects
are cofibrant. The change-of-enrichment theorem therefore makes
restriction a Quillen equivalence. Applying its derived unit to
an objectwise fibrant replacement of the enriched diagram gives a
zigzag of objectwise weak equivalences between that diagram and the
restriction of a diagram on $R^{\mathrm{op}}$. Restriction to the
discrete hom-sets defining $\mathrm{Ty}$ gives the asserted coherent
descent. This checks the hypotheses of the model-category argument
in the intake, using \cite[Proposition A.3.3.8]{CpathsLurieHTT}.

The coherent-nerve argument has the same comparison: replace every
hom-object by $\operatorname{Sing}BP(\tau,\sigma)$ and every value
by $\operatorname{Sing}BF_\tau$. Realization and singular complexes
preserve finite products in these models. All hom-objects are now
Kan; the functor to $R$ is a Dwyer--Kan equivalence and hence its
coherent nerve is an equivalence. Restricting functors along an
equivalence is an equivalence. This proves exactly the descent used
in PATHS, without requiring a strict choice of contractions.

For $\sigma\in R$, the comma category $(\sigma\downarrow p)$
is the opposite of the full subcategory $\mathrm{Ty}_\sigma$
on the types which admit some map to $\sigma$. There is no specified
choice of such a map in this comma category. The functor
$H(\tau)=\sigma*\tau$ has natural maps to $\tau$ and to the
constant object $\sigma$, obtained by collapsing the other block.
It stays in $\mathrm{Ty}_\sigma$, and contracts that category.
For the sequence, the corresponding comma category is the nonempty
upward-closed subposet of $\mathbb N$ consisting of $n$ with
$t_n\to\sigma$. It has a least object. The comma criterion proves
both finality statements.

The fibred Grothendieck comparison identifies $BE_t$ and
$B\operatorname{Ch}_\Sigma$ with the respective homotopy colimits
on $\mathbb N$ and $\mathrm{Ty}^{\mathrm{op}}$. Its naturality
retains the specified inclusion $E_t\to\operatorname{Ch}_\Sigma$.
One may obtain this variance directly by taking opposite categories
and using the covariant theorem; reversal of nerve simplices gives
the natural identification of the two realizations. Descent and
the two finality calculations identify both homotopy colimits with
the colimit over $R^{\mathrm{op}}$. Thus the specified inclusion is
a weak equivalence. The map $E_t\to C_\infty(t)$ is also one:
the nerve commutes with filtered colimits of categories, since an
$n$-simplex uses finitely many objects, arrows and equations, and
filtered colimits of simplicial sets compute their homotopy colimits.
The latter assertion follows by factoring to a projectively
cofibrant diagram and observing, by finite simplicial sphere and
homotopy representatives, that a filtered colimit preserves
objectwise weak equivalences.

For a strictly commuting $u:t\to s$, the cartesian arrows
$(t_n,u_n^*c)\to(s_n,c)$ give a natural transformation in the total
category. Consequently $\beta_t B(u^*)=\beta_s$. It follows that
the actual induced functor $u^*$ is a weak equivalence. The
concatenated sequence has block projections to both sequences and
is cofinal. Its comparison is $\beta_t^{-1}\beta_s$; this formula
proves independence and transitivity in $\mathrm{Hot}$.
\end{proof}

\begin{claim}{C.4}{Proved}{The actual padding and gluing operations}
For a block $\tau_0$ containing both signs, the left, right and
symmetric systems of VII ms.~125--126 have weakly equivalent
colimits $L,R,P$. The actual constant-padding maps $q_L:L\to P$
and $q_R:R\to P$ are weak equivalences. On each endpoint fibre,
\[
 \mu=B\kappa\bigl((Bq_L)^{-1}\times(Bq_R)^{-1}\bigr)
\]
is associative, unital, and independent of the block under the
specified comparisons in $\mathrm{Hot}$.
\end{claim}
\begin{proof}
The alternating-run lengths of $\tau_0^n$ tend to infinity.
For each target word choose, in order, positions with the required
signs in a sufficiently long power. This is a nondecreasing
sign-preserving letter map from the target letters to the power's
letters. The formula in C.11 turns it into a type map from the
power to the target. Thus C.3 applies to all three systems.
If $d$ is the block length, the transition vertex maps are,
respectively, $i\mapsto\max(0,i-pd)$, $i\mapsto\min(i,nd)$,
and $i\mapsto\min(2nd,\max(0,i-pd))$. The maps inducing $q_L$
and $q_R$ at stage $n$ are $i\mapsto\min(i,nd)$ and
$i\mapsto\max(0,i-nd)$ from $\tau_0^{2n}$ to $\tau_0^n$.
Substitution verifies the two strictly commuting squares with
transitions. C.3 therefore applies to these specified maps.

Let $\gamma$ denote literal concatenation on the all-type spaces.
On the common stage with two words $\tau_0^n$, both sides of
\[
 \beta_P B\kappa=\gamma(\beta_L\times\beta_R)
\]
send the pair to its literal concatenation of type $\tau_0^{2n}$.
The transition pads the two outer ends, exactly as in the symmetric
system. The diagonal $\mathbb N\to\mathbb N^2$ is final, and
finite products of spaces commute with filtered colimits, so this
stage calculation proves the equality of the specified maps.
Together with $\beta_P Bq_L=\beta_L$ and
$\beta_P Bq_R=\beta_R$ it gives
\[
 \beta_P\mu=\gamma(\beta_P\times\beta_P).
\]
The strict associativity and empty-path units for $\gamma$ prove
the assertions. The same identity for any other mixed block proves
independence. It is unnecessary to infer equality of maps in
$\mathrm{Hot}$ from equality on a selected class of test objects.
\end{proof}

\begin{claim}{C.5}{Proved}{Exact one-sign answer at VII ms.~125--127}
For every integer $d\geq1$, the left, right and symmetric colimits
formed with the block $+^d$ are isomorphic in $\mathrm{Cat}_{\mathrm{hot}}$
to $\operatorname{Hom}_M(a,b)_{\mathrm{disc}}$; those formed with
$-^d$ are isomorphic there to
$\operatorname{Hom}_\Sigma(b,a)_{\mathrm{disc}}$.
Both actual padding maps are isomorphisms in
$\mathrm{Cat}_{\mathrm{hot}}$. The operation induced by gluing is
associative, unital, and independent of $d$ within either sign.
It is ordinary composition in $M$ in the positive case and
composition in $\Sigma^{\mathrm{op}}$ in the negative case.

The comparison with the all-type localization mapping space, at
every pair of endpoints, is an equivalence for positive blocks
exactly when $\Sigma\subset\operatorname{Iso}(M)$, and for
negative blocks exactly when $M=\Sigma$ is a groupoid.
\end{claim}
\begin{proof}
For a positive path, compose all its arrows to obtain
$f:a\to b$. A transformation fixing endpoints preserves $f$ by
the successive commuting squares. For a negative path the same
argument preserves the composite $w:b\to a$ in $\Sigma$.
At a fixed nonzero length, the factorization category of $f$ has
an initial object with its entire arrow at the last edge and a
terminal object with its entire arrow at the first edge. The
components of the unique comparison transformations are the prefix
and suffix composites. For a negative word the initial object has
its entire arrow at the first edge and the terminal object at the
last edge. All of its comparison components even belong to $\Sigma$.
This proves, in particular, the component and weak-homotopy claims
at each finite stage.

Here is the stronger assertion about the actual colimits. Identify
the positive symmetric colimit with nondecreasing chains indexed by
$\mathbb Z$, eventually equal to $a$ at the left and $b$ at the
right. Multiples of $d$ are cofinal among finite support bounds,
which gives this identification as an isomorphism of categories.
Let $s(f)$ have its unique possibly nonidentity arrow at $0\to1$.
For a path $c$ define $U(c)(j)=c(j)$ for $j\leq0$, and
$U(c)(j)=b$ for $j\geq1$. At the new junction use the suffix
composite $c(0)\to b$. Prefixes and suffixes give the natural zigzag
\[
 c\longrightarrow U(c)\longleftarrow s(f).
\]
Its endpoint components are identities. Every path transformation
preserves the prefix and suffix composites by its commuting
squares, proving naturality in $c$. Thus the effect map and $s$
are inverse in $\mathrm{Cat}_{\mathrm{hot}}$. For a negative
chain, with arrows $j+1\to j$, the same definition uses the suffix
$b\to c(0)$ in $\Sigma$ and gives
$c\leftarrow U(c)\to s(w)$. For the left half-chain the canonical
single-arrow factorization at $-1,0$ is initial or terminal as
appropriate; for the right half-chain use $0,1$. The prefix/suffix
transformations again give the inverse in
$\mathrm{Cat}_{\mathrm{hot}}$.

Both padding maps preserve effect. Their source and target effect
maps are isomorphisms in $\mathrm{Cat}_{\mathrm{hot}}$, so they
are isomorphisms there too. Concatenation preserves effect and sends
$(f:a\to b,g:b\to c)$ to $gf$, or
$(w:b\to a,v:c\to b)$ to $wv$. This proves all composition and
unit assertions, including compatibility for the actual maps.

Neither sign has cofinal powers: no negative letter maps
sign-preservingly into a positive word, and conversely. Under the
comparison in C.8 the effect maps into the all-type space as
$f\mapsto\ell(f)$ and $w\mapsto\ell(w)^{-1}$, where
$\ell:M\to L(M,\Sigma)$ is localization. For a fixed $(a,b)$,
the exact criterion is that the corresponding mapping space be
homotopy discrete and this map of component sets be bijective.
For all endpoints, the positive criterion says precisely that
$NM\to L(M,\Sigma)$ is fully faithful; it is already essentially
surjective. Thus it is an equivalence, and every marked arrow was
already invertible in $M$. The converse is the localization
universal property.

For the negative criterion, the identity-on-objects functor
$N\Sigma^{\mathrm{op}}\to L(M,\Sigma)$, given by inverse
marked arrows, is an equivalence. Every arrow in its image is an
equivalence, so every arrow of $\Sigma$ is invertible in $\Sigma$.
Then $\Sigma$ is a subcategory of the isomorphisms of $M$ and
$L(M,\Sigma)\simeq NM$. Full faithfulness says that inversion
maps $\operatorname{Hom}_\Sigma(b,a)$ bijectively onto
$\operatorname{Hom}_M(a,b)$. Its image is exactly
$\operatorname{Hom}_\Sigma(a,b)$, so $M=\Sigma$, a groupoid.
The converse follows by taking inverses in that groupoid.
\end{proof}

\begin{claim}{C.6}{Refuted, with replacement C.5}{The crown and the scope of the one-sign failure}
For the four-object poset with nonidentity arrows
$0\to1,0\to3,2\to1,2\to3$, mark all arrows and take $a=b=0$.
Every nonempty one-sign fixed-type path category is a point, whereas
the all-type mapping space has component set $\mathbb Z$.
Thus one-sign blocks do not in general compute the all-type object
or provide the comparison with a mixed block. They nevertheless
have the padding and associative unital composition proved in C.5.
\end{claim}
\begin{proof}
A chain in a poset beginning and ending at the same object is
constant when all its arrows have one direction. There is exactly
one such functor and its only transformation is the identity.
The nerve of this crown is the four-cycle. Its covering poset is
the integer fence with arrows from each even integer to its adjacent
odd integers; reduction modulo four is a categorical covering.
It realizes to the universal covering line. An all-type loop based
at zero lifts uniquely beginning at zero, and its last vertex is
$4k$, $k\in\mathbb Z$. The lift endpoint is invariant under
endpoint-fixed ladders and type changes. Traversing the lifted
line from zero to $4k$ supplies a loop for each $k$. Alternatively,
C.7 identifies the all-type space with $\Omega S^1$, which has
these components. The program \texttt{code/cpaths\_checks.py}
checks the finite crown and the covering rule, as well as every
one-sign path of lengths one through seven for all sixteen wide
markings. The unbounded conclusion follows from the poset argument.
This is the exact correction to an interpretation of intake F1
which would make existence or associativity alone fail for a
one-sign block.
\end{proof}

\begin{claim}{C.7}{Proved}{R7 F2, with both adjunctions checked}
For every small $(M,\Sigma)$ as above and $a,b\in M$, there is an
equivalence
\[
 B\operatorname{Ch}_\Sigma(M;a,b)\simeq
 \operatorname{Map}_{L(M,\Sigma)}(a,b).
\]
For fixed $b$ it is an equivalence of presheaves under
$h_b=\operatorname{Hom}_M(-,b)$, after the explicit presheaf
replacement below. Each single-type map $h_b\to X_b^\tau$
is already an equivalence after localization of presheaves at
$\Sigma$.
\end{claim}
\begin{proof}
Restriction along $\ell^{\mathrm{op}}$ identifies presheaves on
$L=L(M,\Sigma)$ with presheaves on $M$ sending $\Sigma$ to
equivalences. This follows directly by applying the localization
universal property to $\mathcal S$ and taking opposites. Restriction
is fully faithful and has left adjoint left Kan extension. Left Kan
extension carries $h_b$ to the representable at $\ell b$, as can
also be checked by Yoneda against every presheaf on $L$. Hence
$R_b(a)=\operatorname{Map}_L(a,b)$ is characterized as the
$\Sigma$-local presheaf under $h_b$ satisfying
\[
 \operatorname{Map}(R_b,Q)\xrightarrow{\sim}Q(b)
 \quad\text{for every $\Sigma$-local $Q$.}
\]

Let $D_\tau$ be the category of $\tau$-paths ending at $b$, with
free origin; transformations have last component $1_b$ and all
other components arbitrary. The origin functor is
$s_\tau:D_\tau\to M$. Its fibred all-type construction is
$s:D\to M$. Define the strict presheaves
\[
 X_b^\tau(a)=B(a\downarrow s_\tau),\qquad
 X_b(a)=B(a\downarrow s).
\]
Precomposition of $a\to s(c)$ defines their strict variance in $a$.
The comma category is literally the fibred Grothendieck construction
of the single-type comma categories, so the natural Grothendieck
comparison gives
$X_b\simeq\operatorname{colim}_{\mathrm{Ty}^{\mathrm{op}}}
X_b^\tau$ in presheaves of spaces. For fixed $\tau$ the comma
projection is the category of elements of $c\mapsto
\operatorname{Hom}_M(a,s_\tau c)$. Its nerve is the usual bar
construction for that homotopy colimit. Consequently
\[
 X_b^\tau\simeq\operatorname{colim}_{c\in D_\tau}h_{c_0},
 \qquad
 \operatorname{Map}(X_b^\tau,Q)
 \simeq\lim_{D_\tau^{\mathrm{op}}}Q(c_0)=:A_\tau.
\]
All comparisons are natural in $a$, because precomposition commutes
with every bar face and degeneracy. The empty type gives $X_b^\varnothing=h_b$.

Define $i:\operatorname{Ch}_\Sigma(M;a,b)\to a\downarrow s$
by $c\mapsto(c,1_a)$, and define $r$ by prepending the forward
edge $f:a\to c_0$ to $(c,f)$. A comma morphism has the equation
$\eta_0 f=f'$. On that morphism $r$ uses the type map with an
extra initial edge and transformation $(1_a,\eta_0,\ldots,\eta_n)$;
the displayed equation is its first square. Collapsing an initial
identity gives $ri\Rightarrow1$. Collapsing the extra initial
edge with components $(f,1,\ldots,1)$ gives $ir\Rightarrow1$
in the comma category; its origin component satisfies
$f1_a=f$. Thus $Bi$ and $Br$ are inverse equivalences. Under this
identification, precomposition by $f$ is prepending its forward edge,
with the comparison just written.

For a marked $w:a'\to a$, forward prepending has inverse, up to
the following natural transformations, given by backward prepending.
The two composites contain the prefixes
\[
 a\xleftarrow{w}a'\xrightarrow{w}a,
 \qquad
 a'\xrightarrow{w}a\xleftarrow{w}a'.
\]
The first transforms to the constant triple at $a$, and the constant
triple at $a'$ transforms to the second. In both cases the middle
component is $w$ and the endpoints are identities. Each square is
an identity equation with $w$, and these transformations are
natural in the remaining path. Collapsing the constant edges gives
the two inverse homotopies. Therefore $X_b$ is $\Sigma$-local.

Now let $Q$ be $\Sigma$-local. Prepending an identity edge defines
$\mathrm{pad}:D_{\tau'}\to D_{\varepsilon\tau'}$ and deleting
the first vertex defines $\mathrm{trunc}$ in the reverse direction.
For $\varepsilon=+$, a map from
$(c_0\xrightarrow{f}c_1;c')$ to $\mathrm{pad}(d')$ is uniquely
determined by a map $c'\to d'$: its first component must be
$\eta_1 f$. Hence $\mathrm{trunc}\dashv\mathrm{pad}$.
The functor $\mathrm{pad}^{\mathrm{op}}$ is a left adjoint and
is initial for limits. Restriction therefore gives an equivalence
$A_{+\tau'}\to A_{\tau'}$ for every $Q$. This also verifies the
intake's opfibration calculation: its fibre is $M/c'_0$, and the
opposite has initial object $1_{c'_0}$, whose value is $Q(c'_0)$.

For $\varepsilon=-$, a map from $\mathrm{pad}(c')$ to
$(d_0\xleftarrow{w}d_1;d')$ is uniquely determined by
$\eta':c'\to d'$, with first component $w\eta_1$.
Thus $\mathrm{pad}\dashv\mathrm{trunc}$, so
$\mathrm{trunc}^{\mathrm{op}}$ is initial. The arrows $w$ give
a natural equivalence from the coefficient diagram $Q(c_0)$ to
$Q(c_1)=Q(s_{\tau'}\mathrm{trunc}(c))$, precisely because $Q$
inverts $\Sigma$. After restriction along $\mathrm{pad}$ this
equivalence is the identity. Therefore restriction
$A_{-\tau'}\to A_{\tau'}$ is again an equivalence.

Every unique type map $\tau\to\varnothing$ factors into these
initial-edge collapses. The diagram $A:\mathrm{Ty}\to\mathcal S$
thus has an objectwise equivalence to its constant value
$A_\varnothing=Q(b)$. The equality of its composites with the
terminal arrows shows, in particular, that it sends all type maps
to equivalences. Since $\mathrm{Ty}$ has terminal empty type,
$B\mathrm{Ty}$ is contractible. It follows that
\[
 \operatorname{Map}(X_b,Q)
 \simeq\lim_{\mathrm{Ty}}A\simeq Q(b).
\]
This is the required localization property of $h_b\to X_b$.
The same calculation at a single $\tau$ proves
$\operatorname{Map}(X_b^\tau,Q)\simeq Q(b)$ and the claimed
single-type local equivalence. It does not assert that a fixed
$X_b^\tau$ itself inverts $\Sigma$.
\end{proof}

\begin{claim}{C.8}{Proved}{The specified map $Q$, both endpoint variables, and composition}
The canonical map $Q_{a,b}$ of PATHS B.6 is the equivalence in C.7.
These maps are natural in $a\in M^{\mathrm{op}}$ and $b\in M$,
with coherent endpoint actions, and preserve composition, units,
and all their coherences. They identify the operation of C.4 with
composition in $L(M,\Sigma)$.
\end{claim}
\begin{proof}
Let $p_\tau$ be the number of positive letters, and let
$q_\tau:I_\tau\to[p_\tau]$ count the positive letters preceding
a vertex. This localizes the interval at its negative edges:
inverting those edges identifies each negative run to a vertex,
leaving the free linear category on the positive edges. This is
also a proof of the universal property in $\infty$-categories:
in a functor to an $\infty$-category, the data of a finite string of
equivalences with fixed value at one end have contractible space of
choices; deleting each such string gives the stated equivalence
of functor categories, including natural transformations.
Thus, for $L=L(M,\Sigma)$, pullback is an equivalence
\[
 \operatorname{Fun}([p_\tau],L)\simeq
 \operatorname{Fun}_\Sigma(I_\tau,L),
\]
where the right side is the full subcategory sending negative edges
to equivalences. It commutes with evaluation at both endpoints.
Take the $\infty$-categorical fibres over $(a,b)$. The endpoint
arrow $[1]\to[p_\tau]$ then gives a functor from that fibre to
$\operatorname{Map}_L(a,b)$, a space. The case $p_\tau=0$ uses
the constant functor $[1]\to[0]$, and includes empty paths.

Every type map $u$ gives a unique monotone $\overline u$ with
$q_\sigma u=\overline u q_\tau$. Therefore these equivalences
form an equivalence of entire diagrams on $\mathrm{Ty}^{\mathrm{op}}$.
Invert this equivalence in that functor $\infty$-category, retaining
its unit and counit; the space of such inverse data is contractible.
The maps from $NF_\tau(a,b)$ to the relative-diagram fibres and
then to $\operatorname{Map}_L(a,b)$ consequently form a coherent
cocone. Taking groupoid completions and the colimit gives exactly
$Q_{a,b}$, through the natural Grothendieck comparison.

For a comma object $(c,f:a\to c_0)$ the same construction evaluates
the path after $\ell(f)$. The comma equation makes this a map of
strict presheaves of spaces $X_b\to R_b$, natural in $a$ and
under $h_b$. By C.7 both its source and target are localizations of
$h_b$. Maps between such localizations under $h_b$ form a
contractible space and are equivalences. This proves that the
specified cocone, rather than an unspecified isomorphism of
homotopy types, is an equivalence.

To check every composition coherence simultaneously, use a finite
list of intervals with marked junctions. Localize each block as
above and retain all the junction evaluations. The functor from
the ordinal of junctions to the concatenated localized ordinal
commutes with erasing junctions and with inserting an empty block.
These are respectively the face and degeneracy operations.
Inverting the single equivalence of diagrams on all such lists
therefore gives a map of Segal objects in spaces. Its source in
degree $n$ is
\[
 \coprod_{a_0,\ldots,a_n}\prod_{j=1}^{n}
 B\operatorname{Ch}_\Sigma(M;a_{j-1},a_j),
\]
with faces given by literal concatenation and degeneracies by empty
paths. Its target is the Segal object for $L$ with this chosen
discrete object set: take the core of $\operatorname{Fun}([n],L)$
and pull back its vertex evaluations along the set of the chosen
objects. Degree one is the map $Q$. This construction verifies
composition and all higher associativities and units at once.
Prepending a forward edge gives the action of $M^{\mathrm{op}}$;
appending one gives the action of $M$. Erasing the junction between
two such edges gives their composite action, and inserting an
empty block gives the unit coherence. The empty path at $b$
evaluates to $1_b$. Finally the identity in C.4 transports this
same composition to the padded systems.
\end{proof}

\begin{claim}{C.9}{Proved}{F3 and the restricted-vertical inclusion}
For every wide $\Sigma$, the specified inclusion
$J:\operatorname{Ch}_\Sigma^\Sigma(M;a,b)\to
\operatorname{Ch}_\Sigma(M;a,b)$ is a weak equivalence.
The identification with Dwyer--Kan hammock localization is natural
in both endpoints, under the ordinary arrows and the empty paths,
and compatible with concatenation.
\end{claim}
\begin{proof}
Write $D_\tau=NF_\tau^\Sigma(a,b)$. A simplex of the restricted
total category is a composable string
$(\tau_0,c_0)\to\cdots\to(\tau_n,c_n)$. Pull every row back
along the composite type map to $\tau_0$; the transformations form
a width-$n$ hammock. Reducing identity columns and composing
same-direction columns gives an $n$-simplex of
$L^H(M,\Sigma)(a,b)$. The zeroth face changes the common first
type, and the two choices agree by the defining colimit relation;
the other faces compose or omit a row, and degeneracies insert an
identity row. This is an actual simplicial map
\[
 \rho:N\operatorname{Ch}_\Sigma^\Sigma(M;a,b)
 \longrightarrow\operatorname{colim}_{\mathrm{Ty}^{\mathrm{op}}}D_\tau
 =L^H(M,\Sigma)(a,b).
\]
Pullback and reduction commute with concatenation, prepending,
appending, and the empty path. In particular these assertions
concern this map, with its structural map from $M$ retained.

Here is the homotopy-colimit check for $\rho$. Give
$\mathrm{Ty}$ degree equal to length, direct maps injective on
vertices and inverse maps surjective on vertices. Factoring through
the image on vertices is the unique Reedy factorization. In the
opposite category the latching map for $D$ is the union of the
subsets in which a nonempty prescribed set of horizontal edges is
identity. Each is a simplicial subset, and intersections correspond
to collapsing the union of the prescribed edges; hence the
latching maps are monomorphisms. The diagram is Reedy cofibrant.
The opposite Reedy category has fibrant constants: the category
of proper injective endpoint-preserving type maps into a given
interval is empty for a reduced interval, and otherwise has initial
object the inclusion of the endpoints and turning vertices. Every
such map contains those vertices, because a comparable pair cannot
cross a missing turn. Thus its nerve is connected. The Reedy
colimit is a left Quillen functor and its ordinary colimit computes
the homotopy colimit. The natural Grothendieck comparison identifies
the induced map with $\rho$, by its first-type formula. This
proves $\rho$ is a weak equivalence and verifies the argument of
Arakawa--Cnossen, Proposition~4.8 and Lemma~4.16
\cite{CpathsArakawaCnossen} in the exact variance used here.

For completeness the marking convention in that paper requires
all equivalences. Put $\Sigma^+=\langle\Sigma,\operatorname{Iso}(M)\rangle$.
The two markings have identical ordinary saturation, since every
functor already preserves isomorphisms. Dwyer--Kan's closure lemma
\cite[Lemma~5.1]{CpathsDKSimplicial}, together with the natural
comparison in \cite[Proposition~2.2]{CpathsDKCalculating}, implies
that $L^H(M,\Sigma)\to L^H(M,\Sigma^+)$ is an equivalence on
mapping spaces. The intermediate comparison diagrams are natural
in the marking, so this conclusion applies to the actual inclusion.
Now Arakawa--Cnossen Theorem~4.7 and Corollary~1.11 apply to
$\Sigma^+$ and identify its mapping spaces with those of
$L(M,\Sigma^+)$. The latter is canonically $L(M,\Sigma)$:
their universal properties select the same functors into every
$\infty$-category. All maps preserve the structural functor from
$M$ and concatenation. This supplies the exact unrestricted-wide
scope required in F3.

The two cocones $Q\,BJ$ and the canonical hammock comparison
evaluate each restricted ladder by inverting its backward edges
and composing its forward edges. They agree as cocones, including
type changes, by the interval-localization construction of C.8.
Thus $Q\,BJ$ is the specified equivalence represented by $\rho$.
Since $Q$ is an equivalence, $BJ$ is an equivalence.

This also supplies a rigorous presheaf replacement for the
coherent restricted path action in F3. Take an enriched fibrant
replacement of $L^H$ and its strict mapping-space presheaf on $M$
using the structural functor $M\to L^H$. The map $\rho$ is an
objectwise equivalence, commutes with prepending followed by
reduction, and preserves the one-edge and empty representatives.
The type-collapse maps between a concatenated pair of forward
edges and its composite map to the identity after reduction;
the same holds for identity insertion. On arbitrary finite lists
of such edges these equalities hold simultaneously by composition
of the type maps. They are exactly the composition and unit
coherences of the action. Therefore this strict presheaf rectifies
the action under $h_b$. The proof of $J$ above uses the specified
cocone square and does not require silently treating the original
total path assignment as a strict presheaf.
\end{proof}

\begin{claim}{C.10}{Proved}{Finite verification of the stated formulas}
The attached independent program verifies the two descriptions of
type maps on all $961$ pairs of words of length at most four, their
contravariant composition on $2731$ composable pairs of maps with
word lengths at most three, the run-clamping formulas on $1213$
nonempty mapping posets, the one-sign crown calculation for $224$
marking/length/sign choices, and $16$ strictly commuting padding
squares for the block $(+,-,+)$. It also verifies the integral
certificate used in C.11 on $7937$ arrows to $1023$ target words.
\end{claim}
\begin{proof}
The program enumerates nondecreasing endpoint-preserving vertex
functions and tests the image of every generating edge in the
oriented target interval. Independently it enumerates nondecreasing
sign-preserving letter maps. It compares the explicit bijections
and checks composition as tuple composition. For the clamp it tests
well-definedness, idempotence, the pointwise natural transformation,
preservation of pointwise order, and the exact image identification
with the smaller-run problem. For padding it substitutes the
vertex functions of C.4 into both sides of each actual square.
The output is \texttt{code/cpaths\_checks.out}; the source contains
the bounds. These are finite audits of the formulas. The proofs in
C.2--C.9 and C.11 supply their unbounded quantifiers.
\end{proof}

\begin{claim}{C.11}{Proved}{R8 with an explicit integral obstruction}
The category $\mathrm{Ty}^{\mathrm{op}}$ is isomorphic to the
category of finite sign words with nondecreasing sign-preserving
letter maps. For $p=(++)$, $q=(+-+)$, the two arrows
$e=(1,3),c=(1,1):p\rightrightarrows q$ have no completion to a
commutative square. They have no affine equalizer over any nonzero
unital ring. The covariant set-valued quotient
\[
 H=\mathrm{Ty}^{\mathrm{op}}(p,-)/(ge\sim gc)
\]
has $H_1(B\operatorname{El}(H);\mathbb Z)=\mathbb Z$.
In particular the opposite type category is not totally
$W_\infty$-aspherical, and categorical colimits over it do not
compute homotopy colimits for every diagram, even for diagrams of
sets.
\end{claim}
\begin{proof}
For a vertex map $\phi:I_\tau\to I_\rho$, associate to the
$j$th target letter the unique source letter $i$ such that
$\phi(i-1)<j\leq\phi(i)$. This is sign-preserving and
nondecreasing. Conversely a letter map $f$ gives
$\phi_f(i)=|\{j:f(j)\leq i\}|$. These formulas are inverse,
including for empty words, and their defining inequalities give
$f_{\psi\phi}=f_\phi f_\psi$. They prove the claimed categorical
isomorphism with its variance.

For every $g:q\to\sigma$ the alternating signs imply
$g(1)<g(2)<g(3)$. Thus $ge$ is injective on the two letters,
whereas $hc$ is constant for every $h:q\to\sigma$. Equality is
impossible, with no bound on $\sigma$.

For a unital ring $R$, define a functional on the free $R$-module
on maps $p\to\sigma$ by assigning value one to injective maps
and zero to constant maps. Then
\[
 \lambda_\sigma\bigl(t(e-c)\bigr)=\operatorname{aug}(t)
 \quad\text{for }t\in R[\mathrm{Ty}^{\mathrm{op}}(q,\sigma)].
\]
An equalizing affine combination would have both sides zero and
augmentation one. This is impossible when $R$ is nonzero.

Over $\mathbb Z$, put $P_x=\mathbb Z[\mathrm{Ty}^{\mathrm{op}}(x,-)]$
and $d:P_q\to P_p$, $d(t)=t(e-c)$. Write $K=\ker d$ and
$J_0=\operatorname{im}d$. Representable free abelian diagrams are
projective and have colimit $\mathbb Z$ by augmentation.
The functional formula makes the map
$\operatorname{colim}K\to\operatorname{colim}P_q$ zero,
so right exactness applied to $K\to P_q\to J_0\to0$ gives
$\operatorname{colim}J_0=\mathbb Z$. Its map to
$\operatorname{colim}P_p$ is zero, since every $t(e-c)$ has
augmentation zero. Free abelian groups preserve the set coequalizer,
so $\operatorname{coker}d=\mathbb Z[H]$. The derived-colimit
sequence of $0\to J_0\to P_p\to\mathbb Z[H]\to0$ therefore
gives $L_1\operatorname{colim}\mathbb Z[H]=\mathbb Z$.
The bar resolution of the permutation diagram $\mathbb Z[H]$
is the nerve-chain complex of its category of elements, proving
the asserted homology. This directly verifies the stronger claim
in TOTAL\_EXACT\_2, C.40, rather than relying only on the span
criterion. The category of elements is the homotopy colimit of the
set diagram, while its categorical colimit is discrete; their first
homology differs.

Here is also a direct check of total asphericity in its final-functor
definition. Put $A=\mathrm{Ty}^{\mathrm{op}}$ and adjoin an object
$*$ with $\operatorname{Hom}(*,a)=H(a)$, no arrows $a\to *$,
and the compositions prescribed by $H$; call the resulting category
$A_H$. The set $\operatorname{colim}H$ is a singleton: it is a
quotient of the colimit of the representable $A(p,-)$, which is a
singleton. Hence $\operatorname{El}(H)$ is nonempty and connected.
The inclusion $A\to A_H$ is an ordinary final functor: its comma
category at $*$ is $\operatorname{El}(H)$, and its comma category
at an old object $a$ is $a\downarrow A$, with initial object $1_a$.
The nerve of $A_H$ is obtained from $NA$ by adjoining the cone on
$N\operatorname{El}(H)$ along the projection. Indeed a simplex
not in $NA$ starts at $*$ and is exactly a chosen element followed
by a simplex of $\operatorname{El}(H)$; this description respects
all faces and degeneracies. Thus $BA_H$ is the homotopy cofiber of
$B\operatorname{El}(H)\to BA$. The empty word is initial in $A$,
so $BA$ is contractible, and the cofiber sequence gives
$H_2(BA_H;\mathbb Z)=H_1(B\operatorname{El}(H);\mathbb Z)=\mathbb Z$.
An ordinary final image of the contractible category $A$ has
therefore failed to be weakly contractible. This proves the asserted
failure of total asphericity directly and confirms the application
of TOTASPH C.5 and C.11 used in R8.
\end{proof}

\begin{claim}{C.12}{Verified derivations}{Literature scope of the path comparison}
The all-type equivalence of C.7--C.9 follows from the proofs above
and the specified hammock comparison, with no calculus-of-fractions
hypothesis. The primary sources opened for this section support the
comparisons summarized below. No priority is claimed.
\end{claim}
\begin{proof}[Source-by-source derivation]
Dwyer--Kan's \emph{Simplicial localizations} \cite{CpathsDKSimplicial} defines the standard
localization for a wide subcategory in Definition~4.1; its
Lemma~5.1 permits saturation, and \S5.5 identifies the
invert-everything mapping spaces with loop spaces.
\emph{Calculating simplicial localizations} \cite{CpathsDKCalculating}, Definition~2.1,
requires vertical arrows in the marking; Proposition~2.2 compares
that model with the standard one; Proposition~5.5 is the type-colimit
formula used for $\rho$; Proposition~6.2 supplies the specified
short-fraction comparisons under the stated homotopy calculus.
Thus the map $J$ of C.9 is the precise additional comparison from
these restricted ladders to the arbitrary-vertical model.
\emph{Function complexes} \cite{CpathsDKFunction}, Propositions~4.4--4.5 and
Corollary~4.7, identify the restricted hammock mapping spaces with
resolution complexes and with a simplicial model category's mapping
space at cofibrant--fibrant pairs. Composing those maps with the
inverse of $BJ$ in $\mathrm{Hot}$ gives the corresponding
arbitrary-vertical comparison under exactly those hypotheses.

Low--Mazel-Gee, \emph{From fractions to complete Segal spaces} \cite{CpathsLowMazelGee},
uses marked vertical arrows in its zigzag categories. Its
three-arrow theorem and its Segal-space theorem apply under the
homotopical three-arrow calculus; the complete-Segal conclusion
also has the stated saturation condition. Composing its reduction
map with C.9 gives its consequence for the present all-type model.
Mazel-Gee, \emph{Hammocks and fractions in relative
$\infty$-categories} \cite{CpathsMazelGee}, defines the enriched hammock construction
with relative transformations and proves the short description
under a homotopical three-arrow calculus. The general localization
comparison used here is supplied in the opened Arakawa--Cnossen
source: Definition~1.7, Theorem~1.9, Corollary~1.11,
Theorem~4.7, Proposition~4.8, Lemma~4.16 and the
Grothendieck-construction remark. Its marking convention is handled
explicitly in C.9. These exact statements use restricted vertical
arrows; C.9 gives the morphism transferring them to the specified
unrestricted model.

Cisinski, \emph{Higher categories and homotopical algebra} \cite{CpathsCisinski},
Theorem~7.2.8, computes derived mapping spaces from a chosen
calculus of fractions, and Remark~7.2.10 identifies the associated
span realization. Carranza--Kapulkin--Lindsey \cite{CpathsCarranzaKapulkinLindsey}, Corollary~10.10,
computes mapping spaces under their quasicategorical calculus of
fractions and its stated marking hypotheses. Arakawa--Carranza--Kapulkin \cite{CpathsArakawaCarranzaKapulkin},
Definition~2.1 and Theorem~2.2, formulates a resolution criterion:
the indexing category is weakly contractible, the structure arrows
become equivalences, and the colimit of representables is local.
These give short or resolution-based models when their hypotheses
are verified; the full path model is verified directly by C.7.

The search also opened Guill\'en--Navarro--Pascual--Roig,
\emph{A Cartan--Eilenberg approach to homotopical algebra},
\S1.2 \cite{CpathsGuillenEtAl}. Its arbitrary-vertical ladders present the ordinary
localization by a quotient of zigzags, which is the $\pi_0$
consequence of C.7. The 2020 MathOverflow question and 2024 answer
on arbitrary vertical arrows \cite{CpathsMathOverflow} were checked as a record of the
question; no theorem in that discussion is used in the proof.
The bibliography records both the supplied source versions and the
web searches. The supplied Arakawa--Cnossen LaTeX was opened in
full at all cited definitions and proofs; two web fetches of that
preprint were disabled, so verification rests on that supplied
source, not an inaccessible web body.
\end{proof}
